\section*{Chapitre \ref{ch:g11:ocean-ridges} --- Les dorsales océaniques et la lithosphère océanique}

\begin{solution}{exo:g11:ocean-ridges:1}
La \omterm{def:g11:ocean-ridges:peridotite}{péridotite} est une roche sombre, à gros grains, faite surtout d'olivine,
avec un peu de pyroxène~; une grande partie du manteau en est faite. Le
\omterm{def:g11:ocean-ridges:gabbro}{gabbro} est une roche magmatique sombre, à gros grains, faite de feldspath
plagioclase et de pyroxène, la roche la plus abondante de la \omterm{def:g11:ocean-ridges:oceanic-lithosphere}{croûte
océanique} profonde. Le \omterm{def:g11:ocean-ridges:gabbro}{gabbro} appartient à la croûte, la \omterm{def:g11:ocean-ridges:peridotite}{péridotite} au
manteau. Le \omterm{def:g11:ocean-ridges:gabbro}{gabbro} a cristallisé lentement, profondément sous le fond marin,
là où le magma refroidissait lentement~: ses cristaux ont eu le temps de
grandir. La lave basaltique a refroidi vite sur le fond marin, et ses
cristaux sont restés petits.
\end{solution}

\begin{solution}{exo:g11:ocean-ridges:2}
Depuis le fond marin vers le bas~: les sédiments (absents à l'axe), les
laves basaltiques (coussins et coulées en nappes), le complexe filonien, le
\omterm{def:g11:ocean-ridges:gabbro}{gabbro}, puis, sous le \omterm{def:g11:seismic-structure:moho}{Moho}, la \omterm{def:g11:ocean-ridges:peridotite}{péridotite} du manteau. Le forage a traversé
$1507/7100 \approx 0.21$, environ un cinquième de la croûte~: aucun forage
n'a jamais atteint les \omterm{def:g11:ocean-ridges:gabbro}{gabbros} profonds ni le \omterm{def:g11:seismic-structure:moho}{Moho}.
\end{solution}

\begin{solution}{exo:g11:ocean-ridges:3}
$d(0) = \qty{2500}{m}$~; $d(9) = 2500 + 350 \times 3 = \qty{3550}{m}$~;
$d(25) = 2500 + 350 \times 5 = \qty{4250}{m}$.
\end{solution}

\begin{solution}{exo:g11:ocean-ridges:4}
Une \omterm{def:g11:ocean-ridges:normal-fault}{faille normale} est une faille inclinée le long de laquelle le bloc situé
au-dessus du plan de faille s'est abaissé par rapport au bloc situé
au-dessous. À une dorsale lente, les plaques écartent la \omterm{def:g11:lithosphere-asthenosphere:lithosphere}{lithosphère} froide
et cassante~: elle est étirée, et elle se rompt. Les blocs situés au-dessus
des failles glissent vers le bas, vers l'axe, et la vallée axiale est le
bloc qui s'est effondré entre eux.
\end{solution}

\begin{solution}{exo:g11:ocean-ridges:5}
Elle ne bout pas à cause de la pression~: plus de deux kilomètres d'eau
pèsent sur elle, et sous cette pression l'eau reste liquide bien au-delà de
\qty{100}{\celsius}. La «~fumée~» noire est un nuage de grains minuscules
de sulfures de fer et d'autres minéraux, qui précipitent quand l'eau chaude
de la source, riche en métaux dissous, rencontre l'eau de mer froide.
\end{solution}

\begin{solution}{exo:g11:ocean-ridges:6}
(1) Vers \qty{54}{km} de profondeur, à environ \qty{1324}{\celsius}.
(2) Vers \qty{85}{km} de profondeur, à environ $1400 + 0.45 \times 85 \approx
\qty{1438}{\celsius}$. (3) Le manteau plus chaud commence à fondre plus
bas~: une plus longue colonne de roche qui remonte fond donc avant
d'atteindre la croûte. Il forme davantage de basalte, et une croûte plus
épaisse.
\end{solution}

\begin{solution}{exo:g11:ocean-ridges:7}
(1) $1489 - 1345 = \qty{144}{\celsius}$ sous sa température de fusion.
(2) Elle remonte de $100 - 54 = \qty{46}{km}$~; l'écart se réduit de
$144/46 \approx \qty{3.1}{\celsius}$ par kilomètre en moyenne. (3) Non~:
pendant la remontée, la roche se refroidit même un peu, d'environ
$0.45 \times 46 \approx
\qty{21}{\celsius}$. C'est sa température de fusion qui baisse, bien plus
vite, à mesure que la pression diminue.
\end{solution}

\begin{solution}{exo:g11:ocean-ridges:8}
Les demi-vitesses valent \qty{7.3}{cm} et \qty{1.65}{cm} par an, soit
\qty{73}{km} et \qty{16.5}{km} par million d'années. Pour former
\qty{1000}{km} de fond d'un côté~: $1000/73 \approx 13.7$ millions d'années
pour la dorsale du Pacifique Est, $1000/16.5 \approx 61$ millions d'années
pour la dorsale médio-atlantique.
\end{solution}

\begin{solution}{exo:g11:ocean-ridges:9}
(1) \qty{5}{cm} valent \qty{5e-5}{km}~: $65000 \times 5 \times 10^{-5} =
\qty{3.25}{km^2}$ de nouveau fond par an. (2) $3.25 \times 7.1 \approx
\qty{23}{km^3}$ de croûte par an. (3) $3.6 \times 10^8 / 3.25 \approx 1.1
\times 10^8$ ans, environ 110 millions d'années~: à ce rythme, les dorsales
pourraient former tout le fond océanique en une centaine de millions
d'années, une durée courte au regard de l'âge de la Terre. C'est pourquoi
aucun fond océanique n'est très ancien.
\end{solution}

\begin{solution}{exo:g11:ocean-ridges:10}
(1) Le plan de faille fait \ang{60} avec l'horizontale, et l'effondrement est
le côté vertical d'un triangle rectangle dont l'hypoténuse est le rejet le
long du plan~: $1.5/\sin\ang{60} \approx \qty{1.73}{km}$. (2) Chaque faille
écarte les bords horizontalement de $1.5/\tan\ang{60} \approx \qty{0.87}{km}$~; les deux failles ensemble, d'environ \qty{1.73}{km}.
(3) Non. Les \qty{25}{km} sont la distance entre les failles qui limitent
le bloc effondré, et non la valeur de l'étirement, qui n'est que d'environ
\qty{1.7}{km}~; et la dorsale continue de s'ouvrir par de nouveaux dykes et
de nouvelles laves à l'axe, en plus du jeu des failles.
\end{solution}

\begin{solution}{exo:g11:ocean-ridges:11}
$d(25) = 6400 - 3200\,e^{-25/62.8} \approx 6400 - 3200 \times 0.672 \approx
\qty{4251}{m}$~; $d(100) = 6400 - 3200\,e^{-100/62.8} \approx 6400 - 3200
\times 0.203 \approx \qty{5749}{m}$. Entre 25 et 100 millions d'années, le
fond s'enfonce d'environ \qty{1500}{m}. La loi en racine carrée donne
$2500 + 350 \times 5 = \qty{4250}{m}$ à 25 millions d'années~: les deux
lois concordent là où toutes deux s'appliquent.
\end{solution}

\begin{solution}{exo:g11:ocean-ridges:12}
(1) $\alpha\,\Delta T\,L = 3.2 \times 10^{-5} \times 700 \times 100 =
\qty{2.24}{km}$. (2) $s = 3380 \times 2.24/(3380 - 1025) = 7571/2355 \approx
\qty{3.2}{km}$. (3) La loi exponentielle permet en tout un enfoncement de
$6400 - 2500 =
\qty{3900}{m}$~: le refroidissement et la contraction suffisent à en
expliquer environ les quatre cinquièmes, avec des valeurs simples pour la
baisse de température et pour l'épaisseur.
\end{solution}

\begin{solution}{exo:g11:ocean-ridges:13}
(1) $54 - 7 = \qty{47}{km}$ contre $85 - 7 = \qty{78}{km}$~: la colonne la
plus chaude est $78/47 \approx 1.66$ fois plus longue. (2) $7.1 \times 1.66 \approx
\qty{11.8}{km}$. (3) Les \qty{10.3}{km} mesurés sont plus épais que la
normale, comme prévu, mais moins que l'estimation~: les kilomètres les plus
profonds de la colonne, juste au-dessus du début de la fusion, donnent moins
de liquide que les moins profonds, si bien que l'hypothèse d'une même
quantité de liquide par kilomètre surestime l'épaisseur.
\end{solution}

\begin{solution}{exo:g11:ocean-ridges:14}
(1) La loi donne \qty{3178}{m}, \qty{4160}{m}, \qty{4782}{m} et
\qty{5267}{m}~; les médianes mesurées sont moins profondes de
\qty{64}{m}, \qty{142}{m}, \qty{317}{m} et \qty{619}{m}. (2) L'écart
augmente avec l'âge. (3) La loi donne la profondeur du fond lui-même,
corrigée de ses sédiments~; les cellules mesurées portent leurs sédiments,
plus épais sur les fonds anciens, et comprennent des monts sous-marins et
des plateaux volcaniques, moins profonds que le fond qui les entoure.
\end{solution}

\begin{solution}{exo:g11:ocean-ridges:15}
(1) $p = 1025 \times 9.8 \times 2500 \approx \qty{25.1}{MPa}$, auxquels
s'ajoutent \qty{0.1}{MPa}~: environ \qty{25.2}{MPa}. (2) C'est au-dessus de la pression
critique, \qty{22.064}{MPa}~: l'eau de la source ne peut pas bouillir,
quelle que soit sa température. (3) La pression atteint \qty{22.064}{MPa} à
la profondeur
$(22.064 - 0.1) \times 10^6/(1025 \times 9.8) \approx \qty{2190}{m}$~:
toute source située à plus d'environ \qty{2200}{m} de profondeur est à
l'abri de l'ébullition.
\end{solution}

\begin{solution}{pb:g11:ocean-ridges:1}
\textbf{1.} $14.6/2 = \qty{7.3}{cm}$ par an, soit \qty{73}{km} par million
d'années.

\textbf{2.} À la profondeur de l'axe de la dorsale, $d(0) = \qty{2500}{m}$.

\textbf{3.} Depuis le fond marin vers le bas~: des laves basaltiques
(coussins et coulées en nappes), le complexe filonien et le \omterm{def:g11:ocean-ridges:gabbro}{gabbro}, environ
\qty{7.1}{km} de croûte en tout~; sous le \omterm{def:g11:seismic-structure:moho}{Moho}, la \omterm{def:g11:ocean-ridges:peridotite}{péridotite} du manteau,
jusqu'à la base de la \omterm{def:g11:lithosphere-asthenosphere:lithosphere}{lithosphère}, encore très mince à l'axe. Des sédiments
s'y déposeront plus tard.

\textbf{4.} Quand les plaques s'écartent, l'\omterm{def:g11:lithosphere-asthenosphere:lithosphere}{asthénosphère} remonte sous la
dorsale. Sa température ne baisse que d'environ \qty{0.45}{\celsius} par
kilomètre, alors que sa température de fusion baisse bien plus vite à mesure
que la pression diminue. Vers \qty{54}{km} de profondeur, les deux se
rejoignent~: la roche commence à fondre, et continue de fondre en
remontant, sans qu'aucune chaleur soit apportée (\omterm{def:g11:ocean-ridges:partial-melting}{fusion par décompression}).

\textbf{5.} Chaque million d'années, la dorsale forme $2 \times 73 = \qty{146}{km}$
de fond sur ses deux côtés~; pour chaque kilomètre de dorsale, le volume de
croûte vaut $146 \times 7.1 \times 1 \approx \qty{1040}{km^3}$.

\textbf{6.} $73 \times t$~: \qty{73}{km}, \qty{292}{km}, \qty{1168}{km} et
\qty{4672}{km}.

\textbf{7.} $2000/73 \approx 27$ millions d'années.

\textbf{8.} Les plaques changent de vitesse et de direction en quelques
dizaines de millions d'années, comme le montre le coude de la chaîne
Hawaï--Empereur (\cref{ch:g11:plate-motions})~; la vitesse actuelle n'a pas
de raison de se maintenir pendant 64 millions d'années.

\textbf{9.} Loin de l'axe, la croûte ne se trouve plus au-dessus d'une
lentille de magma~: elle se refroidit, et il n'y a plus assez de chaleur
en dessous pour entretenir la circulation et chauffer l'eau.

\textbf{10.} À l'ouest de la dorsale du Pacifique Est, le morceau voyage sur
la plaque Pacifique. Vers l'ouest, la plaque Pacifique se termine le long de
fosses profondes, où elle s'enfonce sous ses voisines~: des zones de
subduction (\cref{ch:g11:subduction-zones}).

\textbf{11.} $d(1) = \qty{2850}{m}$, $d(4) = 2500 + 700 = \qty{3200}{m}$,
$d(16) = 2500 + 1400 = \qty{3900}{m}$, $d(64) = 2500 + 2800 = \qty{5300}{m}$.

\textbf{12.} $3900 - 2500 = \qty{1400}{m}$ pendant les 16 premiers millions
d'années, et $5300 - 3900 = \qty{1400}{m}$ entre 16 et 64 millions
d'années~: le même enfoncement, en une durée trois fois plus longue. Le fond
s'enfonce de plus en plus lentement en vieillissant.

\textbf{13.} $350\sqrt{t} = 1500$ donne $\sqrt{t} \approx 4.29$, donc
$t \approx 18.4$ millions d'années.

\textbf{14.} $6400 - 3200\,e^{-64/62.8} \approx 6400 - 3200 \times 0.361
\approx \qty{5245}{m}$, environ \qty{55}{m} de moins que la loi en racine
carrée~: les deux lois concordent à environ \qty{1}{\percent} près.

\textbf{15.} La médiane mesurée est moins profonde que la loi d'environ
\qty{650}{m}. La loi décrit le fond corrigé du poids de ses sédiments~; les
profondeurs mesurées comprennent les sédiments, accumulés pendant plusieurs
dizaines de millions d'années sur un tel fond, ainsi que les monts
sous-marins et les plateaux volcaniques, qui rendent tous le fond mesuré
moins profond.

\textbf{16.} $\alpha\,\Delta T = 3.2 \times 10^{-5} \times 700 = 0.0224$~: la
partie mantellique est plus dense de \qty{2.24}{\percent}, soit
$3380 \times 1.0224 \approx \qty{3456}{kg/m^3}$.

\textbf{17.} La partie mantellique est épaisse de $100 - 7.1 = \qty{92.9}{km}$.
Masse volumique moyenne~: $(7.1 \times 2900 + 92.9 \times 3456)/100 \approx (20590 + 321060)/100
\approx \qty{3416}{kg/m^3}$.

\textbf{18.} $3416 > 3380$~: la \omterm{def:g11:lithosphere-asthenosphere:lithosphere}{lithosphère} ancienne, dans son ensemble, est
plus dense que l'\omterm{def:g11:lithosphere-asthenosphere:lithosphere}{asthénosphère} située dessous d'environ \qty{36}{kg/m^3},
soit environ \qty{1}{\percent}. La croûte légère ne suffit plus à compenser
sa partie mantellique froide et dense.

\textbf{19.} Plus dense que l'\omterm{def:g11:lithosphere-asthenosphere:lithosphere}{asthénosphère}, elle peut replonger dans le
manteau à une fosse~: c'est la subduction, le sujet du chapitre suivant.

\textbf{20.} À 64 millions d'années, le morceau de fond se trouve vers
\qty{5300}{m}~: à environ \textbf{5.3~km} de profondeur, \qty{2.8}{km} plus
bas que là où il est né.
\end{solution}
